\documentclass[12pt]{article}

% --- Packages ---
\usepackage[margin=1in]{geometry}
\usepackage{booktabs}
\usepackage{graphicx}
\usepackage{amsmath, amssymb}
\usepackage{natbib}
\usepackage{hyperref}
\usepackage{setspace}
\usepackage{caption}

\bibliographystyle{chicago}   % Swap for aer / qje / restud as journal-specific files are dropped in
\setcitestyle{authoryear,round}
\doublespacing

% --- Title block ---
\title{\textbf{<PAPER TITLE>}\thanks{Acknowledgments.}}
\author{<Author Name>\\<Affiliation>\\\texttt{<email>}}
\date{\today}

\begin{document}
\maketitle

% =====================================================
\begin{abstract}
\noindent
% 200 words. Lead with the headline number.
% One-sentence question. One-sentence data + identification.
% Two-sentence main result. One-sentence contribution.
\end{abstract}

\medskip
\noindent\textbf{JEL:} % e.g. J21, J31, O15
\\\textbf{Keywords:} % 3-5 keywords

% =====================================================
\section{Introduction}
\label{sec:intro}
% Paragraph 1: problem statement.
% Paragraph 2: why it matters (theory/policy/literature stakes).
% Paragraph 3: what I do (data, identifying variation, method — one sentence each).
% Paragraph 4: main finding (headline number, units, interpretation, brief robustness mention).
% Paragraph 5: contribution (3 bullets vs prior literature) and roadmap.

% =====================================================
\section{Literature}
\label{sec:lit}
% Organize by sub-question, not chronologically.
% Anchor each paragraph on a gap-analysis tension from annotated_bib.md.
% \citet{key} when author is subject; \citep{key} parenthetical.

% =====================================================
\section{Data}
\label{sec:data}
% Source, sample period, unit of observation.
% Sample-construction filters (point to design.md for full detail).
% Key variables with construction notes.
% Summary stats:
% \input{../results/tab_summary.tex}

% =====================================================
\section{Empirical Strategy}
\label{sec:strategy}
% Estimating equation.
\begin{equation}
Y_{it} = \alpha_i + \gamma_t + \beta \, D_{it} + X_{it}' \theta + \varepsilon_{it}
\label{eq:main}
\end{equation}
% Identification argument. Threats. Mitigations. Cluster level.

% =====================================================
\section{Results}
\label{sec:results}

\subsection{Main result}
% \input{../results/tab_main.tex}
% Two-three paragraphs interpreting the coefficient in plain English
% (drawn from results/interpretation.md).

\subsection{Heterogeneity}
% \input{../results/tab_het.tex}
% \begin{figure}[h] \centering \includegraphics[width=.8\textwidth]{../results/fig_het.pdf}
% \caption{Heterogeneous effects.} \label{fig:het} \end{figure}

% =====================================================
\section{Robustness}
\label{sec:robust}
% Walk through each robustness table in 1-2 sentences.
% \input{../results/tab_robust_1.tex}

% =====================================================
\section{Conclusion}
\label{sec:conclusion}
% Recap finding. Situate against literature. 1-2 honest limitations + future directions.

% =====================================================
\section*{Acknowledgments}
% Stub.

\section*{Data availability}
% Describe data source, access, replication code repository.

% =====================================================
\bibliography{refs}

\end{document}
